
% ============================================================
\section{Results \& Conclusion}
% ============================================================

\subsection{Evaluation Results}

\subsubsection{Test Suite}

The 27-module suite is the primary quantitative instrument.
Table~\ref{tab:tests} groups the modules by metric family and
reports the aggregate outcome of the current suite.

\begin{table}[t]
\centering\small
\caption{Automated test suite by metric family (27 modules)}
\label{tab:tests}
\begin{tabular}{@{}p{0.36\columnwidth}c@{\qquad}p{0.2\columnwidth}@{}}
\toprule
\textbf{Family} & \textbf{Modules} & \textbf{Result} \\
\midrule
Safety \& gating (queue, permission, drop, dnd) & 8 & pass \\
Inference (continue, actions, modelprobe, modelradar) & 4 & pass \\
Services (news, digest, brief, weather, ghwatch) & 5 & pass \\
Overlay \& UX (ctrl, fullscreen, region, screenshot) & 5 & pass \\
Ops (config, fsutil, inbox, selfaudit, remind$\ldots$) & 5 & pass \\
\midrule
\textbf{Total} & \textbf{27} & \textbf{all pass} \\
\bottomrule
\end{tabular}
\end{table}

\subsubsection{Failure Injection}

Three injections were run against the assistant path:

\begin{enumerate}
  \item \textbf{Quota exhaustion (tier 1 credit error).} The
  hosted session rejects; the router advances to the next model
  in the chain and the reply completes without operator
  intervention. Subsequent healthy runs confirm no stuck state.
  \item \textbf{Service unreachable (\texttt{ECONNREFUSED}).}
  With the local service port closed, the router skips the
  session tier and answers from the next reachable tier; the
  assistant replies with an explicit offline status instead of
  hanging.
  \item \textbf{Network loss (\texttt{ENOTFOUND}).} With DNS
  resolution failing for all cloud endpoints, Tier~3 (local) or
  Tier~4 (deterministic) returns a substantive message ---
  queue state, cached weather, last digest --- so the user
  always receives an answer. This is the behaviour that
  distinguishes the design from single-provider assistants,
  which surface a raw error.
\end{enumerate}

In every injection the invariant \emph{``no prompt leaves the
system without either a reply or an explicit, actionable status
message''} held. Failover decisions are written to the debug log
with the failing tier and reason, which also feeds the
self-audit reader.

\subsubsection{Operational Reliability}

The watchdog task executes every minute with
\texttt{DisallowStartIfOnBatteries=false},
\texttt{StopIfGoingOnBatteries=false} and
\texttt{StartWhenAvailable=true}; after these settings were
applied, missed-run counts dropped to zero on battery power,
demonstrating that recovery from host-process termination is
bounded by one scheduling interval even after sleep or
shutdown-veto events.

\subsection{Discussion}

Three findings summarize the evaluation.

\begin{enumerate}
  \item \textbf{The four-tier chain preserves continuity where
  single-provider designs fail.} The hosted-session tier
  provided the best latency and answer quality under normal
  conditions; each failure mode --- credit, connection, network
  --- advanced one tier and still produced a usable reply,
  making availability an architectural property rather than an
  operational accident.
  \item \textbf{Bounded FIFO with human gating scales safely to
  phone traffic.} Capacity $Q_{max}=8$ with strict one-at-a-time
  execution gave explicit back-pressure instead of unbounded
  growth, and permission prompts kept every consequential action
  behind a human decision while \emph{always} rules removed
  repeated friction without losing the audit trail.
  \item \textbf{An ambient overlay is a viable control surface
  for a 24/7 agent.} Hover expansion, cursor-away collapse and
  CSS-only visibility control kept the assistant glanceable
  without ever interrupting full-screen work --- including
  automatic dimming during detected full-screen video or games.
\end{enumerate}

\subsection{Conclusion}

Agent Notch demonstrates that an ambient desktop overlay, a
phone-mediated control channel, and provider-agnostic inference
can be composed into one coherent personal-agent system. The
overlay provides Dexter-style glanceable status at the top edge
of the display; WhatsApp provides remote presence with bounded
FIFO execution and human permission gating; and the four-tier
fallback chain --- hosted session, OpenAI-compatible cloud,
local Ollama, deterministic offline --- ensures the assistant
answers under quota exhaustion, service outage and total
network loss. Evaluation through 27 automated test modules and
three failure injections confirms continuous reply availability,
safe one-at-a-time execution and watchdog-bounded recovery.

Future enhancements include widening the OpenAI-compatible tier
with additional provider profiles, adding streaming first-token
metrics to the latency family, and extending the digest engine
with user-defined topic feeds --- all designed to slot into the
existing tier and queue abstractions without changing the
system's safety invariants.

% ============================================================
\begin{thebibliography}{9}

\bibitem{react} S. Yao, J. Zhao, D. Yu, N. Du, I. Shafran,
K. Narasimhan, and J. Cao, ``ReAct: Synergizing reasoning and
acting in language models,'' in \emph{Proc. ICLR}, 2023.

\bibitem{tooluse} T. Schick, J. Dwivedi-Yu, R. Dessi, R. Raileanu,
M. Lomeli, E. Hambro, L. Zettlemoyer, N. Cancedda, and
T. Scialom, ``Toolformer: Language models can teach themselves to
use tools,'' in \emph{Proc. NeurIPS}, 2023.

\bibitem{langchain} H. Chase, ``LangChain: Chainable language
model intelligence,'' \url{https://github.com/langchain-ai/langchain}, 2022.

\bibitem{autogpt} T. Torantoino, ``AutoGPT: Autonomous GPT-4
agent,'' \url{https://github.com/Significant-Gravitas/AutoGPT}, 2023.

\bibitem{ollama} Ollama, ``Run large language models locally,''
\url{https://ollama.com}, 2024.

\bibitem{whatsapp} WhatsApp, ``WhatsApp Business Platform:
Cloud API,'' \url{https://developers.facebook.com/docs/whatsapp}, 2024.

\bibitem{telegram} Telegram, ``Bots: An introduction for
developers,'' \url{https://core.telegram.org/bots}, 2024.

\bibitem{routing} L. Zheng, Z. Liu, G. Lai, and W. Zhang,
``Survey of model routing in large language model systems,''
\emph{arXiv preprint arXiv:2407.14832}, 2024.

\bibitem{ambient} M. Weiser and J. Brown, ``Designing calm
technology,'' \emph{Xerox PARC}, 1995.

\bibitem{hig} Apple Inc., ``Dynamic Island --- Human Interface
Guidelines,'' \url{https://developer.apple.com/design/human-interface-guidelines/dynamic-island}, 2024.

\bibitem{eddata} C. Romero and S. Ventura, ``Educational data
mining and learning analytics: An updated survey,'' \emph{WIREs
Data Mining and Knowledge Discovery}, vol. 10, no. 3, 2020.

\bibitem{dropout} N. R. Lakshmi and M. N. Kalyani, ``Prediction
of at-risk students in e-learning systems: A systematic
literature review,'' \emph{Computers \& Education: Artificial
Intelligence}, vol. 6, 2024.

\bibitem{transformer} A. Vaswani, N. Shazeer, N. Parmar,
J. Uszkoreit, L. Jones, A. N. Gomez, L. Kaiser, and
I. Polosukhin, ``Attention is all you need,'' in \emph{Proc.
NeurIPS}, 2017.

\bibitem{cot} J. Wei, X. Wang, D. Schuurmans, M. Bosma,
B. Ichter, F. Xia, E. Chi, Q. Le, and D. Zhou, ``Chain-of-thought
prompting elicits reasoning in large language models,'' in
\emph{Proc. NeurIPS}, 2022.

\bibitem{elect} Electron Authors, ``Electron: Build cross-platform
desktop apps with JavaScript, HTML, and CSS,''
\url{https://www.electronjs.org}, 2024.

\bibitem{opencode} OpenCode, ``OpenCode agent platform,''
\url{https://opencode.ai}, 2025.

\bibitem{llmsurvey} J. Wu, W. Chen, Y. Su, \emph{et al.},
``Large language model agents: A survey,'' \emph{arXiv preprint
arXiv:2405.04231}, 2024.

\bibitem{rag} P. Lewis, E. Perez, A. Piktus, F. Petroni,
V. Karpukhin, N. Goyal, H. K\"{u}ttler, M. Lewis, W.-t. Yih,
T. Rockt\"{a}schel, S. Riedel, and D. Kiela, ``Retrieval-augmented
generation for knowledge-intensive NLP tasks,'' in \emph{Proc.
NeurIPS}, 2020.

\end{thebibliography}

\end{document}
